\section{Media Transmisi dan Topologi}
\subsection{Media Transmisi}
Sarana atau jalur fisik yang digunakan untuk mengirimkan data/informasi dari sumber pengirim ke perangkat penerima.

\subsection{Klasifikasi Media Transmisi}
$$\text{Media Transmisi}
\begin{cases}
  \text{Kabel (Wired)} \\
  \text{Tanpa Kabel (Wireless)}
\end{cases}
$$

\subsubsection{Kabel (Wired)}
Menggunakan kabel sebagai media penghubung, sinyal dihantarkan melalui media fisik. Biasanya digunakan hanya pada area lokal saja.


\textbf{Contoh media transmisi yang menggunakan kabel:}
\begin{enumerate}
\item Twisted Pair (UTP/STP) \\
  Terdiri atas beberapa pasang kabel tembaga yang dipilin (twisted), menggunakan konektor RJ45.
\item Fiber Optic \\
  Terbuat dari serat kaca, menggunakan cahaya untuk mentransmisikan data. Lapisannya terdiri atas:
  \begin{enumerate}
  \item Core.
  \item Cladding.
  \item Coating.
  \item Strength Member.
  \item Outer Jacket.
  \end{enumerate}
\item Coaxial \\
  Memiliki satu konduktor tembaga di bagian tengah yang dikelilingi lapisan isolator.
\end{enumerate}

\textbf{Karakteristik: }
\begin{table}[H]
  \centering
  \begin{tabular}{cccc}
    \toprule
    \textbf{Jenis} & \textbf{Kecepatan} & \textbf{Jarak} & \textbf{Ketahanan Terhadap Interferensi} \\
    \midrule
    UTP & Rendah hingga 1 gbps & $\pm$ 100 Meter & Rendah  \\
    STP & Rendah hingga 1 gbps & $\pm$ 100 Meter & Sedang - Tinggi  \\
    Coaxial & Sedang hingga 10 gbps & $\pm$ 500 Meter & Sedang  \\
    Fiber Optik & Sangat Tinggi hingga 100 gbps & Puluhan Kilometer & Tinggi \\
    \bottomrule
  \end{tabular}
\end{table}

\subsubsection{Nirkabel (Wireless)}
Teknologi yang memungkinkan transfer data tanpa kabel.
\begin{enumerate}
\item Gelombang Radio \\
  Gelombang yang merambat di udara auat ruang hampa pada kecepatan cahaya dan dimanfaatkan secara luas (omnidirectional).
\item Komunikasi Satelit \\
  Satelit di luar angkasa yang berfungsi sebatai repeater untuk menerima, memperkuat dan memancarkan sinyal untuk dipantulkan kembali ke bumi. Mampu menjangkau wilayah yang sangat luas.
\item Inframerah \\
  Media transmisi jarak dekat, membutuhkan posisi line-of-sight, tidak dapat menembuh dinding atau penghalang padat.
\end{enumerate}

\subsection{Perbandingan Antara Kabel dan Nirkabel}
\subsubsection{Kabel}
\begin{enumerate}
\item Kecepatan tinggi dan stabilitas tinggi.
\item Tahan interferensi eksternal.
\item Instalasi kurang fleksibel.
\item Terbatas oleh panjang kabel.
\end{enumerate}
\subsubsection{Nirkabel}
\begin{enumerate}
\item Fleksibel dan mudah dipasang.
\item Menjangkau wilayah luas dan terpencil.
\item Rentan gangguan dan cuaca.
\item Lebih rawan disadap.
\end{enumerate}
